\section*{Capítulo \ref{ch:b1:weathering-clays} --- Meteorización y arcillas}

\begin{solution}{exo:b1:weathering-clays:1}
\ce{K+}: $1/1.38 = 0.72$; \ce{Fe^2+}: $2/0.78 = 2.56$; ambos por debajo de 3, cationes solubles.
\ce{Fe^3+}: $3/0.645 = 4.65$; \ce{Ti^4+}: $4/0.605 = 6.61$; ambos entre 3 y 10, hidrolizados. En un suelo encharcado, sin oxígeno, el hierro
sale como \ce{Fe^2+}, que permanece en disolución y lo arrastra el agua; en un suelo bien drenado, el oxígeno lo convierte en \ce{Fe^3+}, que
precipita de inmediato como óxidos e hidróxidos y permanece.
\end{solution}

\begin{solution}{exo:b1:weathering-clays:2}
A se hincha con glicol y se colapsa al calentar: una \omterm{def:b1:weathering-clays:clay-mineral}{esmectita}. B está a \SI{7}{\angstrom} y se destruye a \SI{550}{\celsius}: \omterm{def:b1:weathering-clays:clay-mineral}{caolinita}.
C, a \SI{10}{\angstrom} con cualquier tratamiento: \omterm{def:b1:weathering-clays:clay-mineral}{ilita}. D, a \SI{14}{\angstrom} sin hincharse ni colapsar: \omterm{def:b1:weathering-clays:clay-mineral}{clorita}.
\end{solution}

\begin{solution}{exo:b1:weathering-clays:3}
Dos ortoclasas aportan los 2 aluminios de una \omterm{def:b1:weathering-clays:clay-mineral}{caolinita}; $6 - 2 = 4$ silicios salen como \ce{H4SiO4}; salen 2 \ce{K+}, equilibrados por 2 \ce{H+}; el oxígeno,
$16 + w = 9 + 16$, da $w = 9$:
\[
  \ce{2 KAlSi3O8 + 2 H+ + 9 H2O -> Al2Si2O5(OH)4 + 2 K+ + 4 H4SiO4}.
\]
El hidrógeno se comprueba: $2 + 18 = 20 = 4 + 16$. Sale todo el potasio, y $4/6 =
2/3$ del silicio.
\end{solution}

\begin{solution}{exo:b1:weathering-clays:4}
(a) La masa no cambia: $10 \times 999.87/916.72 = \qty{10.91}{cm^3}$.
(b) Entre unos \qty{-4}{\celsius} y \qty{-15}{\celsius}. Justo por debajo de cero el hielo no puede acumular presión suficiente para cuñar la grieta;
con una helada más dura, el agua ya no se desplaza a través de los poros hacia el hielo, que deja de crecer. (c) Por
\cref{prop:b1:weathering-clays:surface} con $n = \qty{1}{m}/\qty{1}{cm} = 100$,
la superficie se multiplica por 100: de \qty{6}{m^2} a \qty{600}{m^2}.
\end{solution}

\begin{solution}{exo:b1:weathering-clays:5}
$S/m = 6/(\rho d)$: para \qty{1}{cm}, $6/(2600\times0.01) =
\qty{0.231}{m^2.kg^{-1}}$; para \qty{1}{mm}, \qty{2.31}{m^2.kg^{-1}}; para
\qty{1}{\micro m}, \qty{2308}{m^2.kg^{-1}}, o \qty{2.3}{m^2.g^{-1}}.
Las reacciones químicas tienen lugar en la superficie de los granos, así que su velocidad por kilogramo de roca crece con la superficie específica,
diez mil veces desde el centímetro hasta el micrómetro. La \omterm{def:b1:weathering-clays:weathering}{meteorización física} corta la roca en trozos cada vez más pequeños y abre grietas por
las que entra el agua: multiplica la superficie sobre la que puede actuar la \omterm{def:b1:weathering-clays:weathering}{meteorización química}.
\end{solution}

\begin{solution}{exo:b1:weathering-clays:6}
Una tonelada de albita son $\num{e6}/262.22 = \qty{3813.6}{mol}$. Según
\cref{prop:b1:weathering-clays:feldspars}, dos moles dan uno de \omterm{def:b1:weathering-clays:clay-mineral}{caolinita}:
$\qty{1906.8}{mol} \times \qty{258.16}{g.mol^{-1}} = \qty{492.3}{kg}$. Salen dos de los tres silicios de cada átomo:
$2\times3813.6 = \qty{7627.2}{mol}$ de \ce{SiO2}, \qty{458.2}{kg}. Sale todo el sodio: $3813.6\times22.99 =
\qty{87.7}{kg}$. Queda como arcilla menos de la mitad de la masa del \omterm{def:b1:rock-forming-minerals:groups}{feldespato}.
\end{solution}

\begin{solution}{exo:b1:weathering-clays:7}
Según \cref{prop:b1:weathering-clays:immobile}, un kilogramo de granito deja
$0.30/0.60 = \qty{0.5}{kg}$ de saprolita. Fracción perdida: sodio $1 -
(0.35/3.5)(0.5) = 0.95$; silicio $1 - (60/72)(0.5) = 0.583$; aluminio $1 -
(27/14)(0.5) = 0.036$. El aluminio se conserva prácticamente, el sodio casi ha desaparecido por completo, y cerca del \qty{60}{\percent} del silicio
también, aunque gran parte de lo que queda está retenido en el cuarzo: los \omterm{def:b1:rock-forming-minerals:groups}{feldespatos} se están convirtiendo en \omterm{def:b1:weathering-clays:clay-mineral}{caolinita}, en una
\omterm{def:b1:weathering-clays:hydrolysis-degree}{monosialitización}.
\end{solution}

\begin{solution}{exo:b1:weathering-clays:8}
Pirofilita: $x = y = 0$, carga 0. Capa de moscovita: $x = 1$, carga $-1$.
Capa de beidellita: $x = 0.5$, carga $-0.5$. Capa de montmorillonita: $y = 0.33$,
carga $-0.33$. Las dos capas de \omterm{def:b1:weathering-clays:clay-mineral}{esmectita}, la de beidellita y la de montmorillonita, pueden
hincharse: su carga pequeña se equilibra con pocos cationes intercapa, que sujetan las
capas con poca fuerza y atraen agua al hidratarse. La pirofilita no tiene cationes intercapa que hidratar, y la carga $-1$ de la moscovita se
equilibra con iones potasio que bloquean las capas.
\end{solution}

\begin{solution}{exo:b1:weathering-clays:9}
(a) Mg: 2 y 2; Si: 1 y 1; C: 4 y 4; O: $4 + 8 + 4 = 16$ y $12 + 4 =
16$; H: 8 y $4 + 4 = 8$; carga: 0 y $4 - 4 = 0$. (b) La masa molar de la
forsterita es $2\times24.305 + 28.085 + 4\times15.999 =
\qty{140.69}{g.mol^{-1}}$, así que una tonelada son \qty{7108}{mol}. Consume
$4\times7108 = \qty{28431}{mol}$ de \ce{CO2}, \qty{1251}{kg}; con dos
magnesios por unidad formular, \cref{prop:b1:weathering-clays:co2} almacena
$2\times7108$ moles, \qty{626}{kg}, la mitad de lo consumido. (c) Una tonelada de
calcita, \qty{9991}{mol}, consume \qty{9991}{mol} de \ce{CO2} (\qty{440}{kg})
al disolverse, y los devuelve todos cuando el carbonato vuelve a precipitar:
no se almacena nada.
\end{solution}

\begin{solution}{exo:b1:weathering-clays:10}
(a) $3 < 3/r < 10$ da $\qty{0.3}{\angstrom} < r < \qty{1.0}{\angstrom}$.
(b) $3/0.27 = 11.1 > 10$: el boro viaja como oxianión, borato. (c) $2/0.83 =
2.41$, un catión soluble; $4/0.53 = 7.55$, un hidrolizado. En agua pobre en
oxígeno el manganeso se disuelve como \ce{Mn^2+}; donde el manantial encuentra el aire,
se oxida a \ce{Mn^4+}, que precipita de inmediato como un óxido negro sobre
la roca. (d) Su \omterm{def:b1:weathering-clays:ionic-potential}{potencial iónico}, $4/0.40 = 10.0$, queda en el límite entre los hidrolizados y los oxianiones: el silicio no se retiene
como el aluminio ni sale como el azufre, y la fracción lixiviada depende de cuánta agua pasa, la base de los tres grados de
\omterm{def:b1:weathering-clays:hydrolysis}{hidrólisis}.
\end{solution}

\begin{solution}{exo:b1:weathering-clays:11}
(a) Con \qty{12.011}{g.mol^{-1}} de carbono, \num{11.7e12} a \qty{13.6e12}{mol} son de 0.141 a 0.163 mil millones de toneladas de carbono al año; con
\qty{44.009}{g.mol^{-1}}, de 0.515 a 0.599 mil millones de toneladas de dióxido de carbono. (b) $4.08/13.6 = \qty{30}{\percent}$ a $4.08/11.7 = \qty{35}{\percent}$.
(c) La mitad del dióxido de carbono consumido se almacena cuando el calcio y el magnesio precipitan como carbonato: de 0.070 a 0.082 mil millones de
toneladas de carbono al año. La cifra real es menor, porque parte del bicarbonato sale con el sodio y el potasio, que no forman carbonato en el fondo marino.
\end{solution}

\begin{solution}{exo:b1:weathering-clays:12}
Si/Al $= (w_{\ce{SiO2}}/60.08)/(2w_{\ce{Al2O3}}/101.96)$. A: $0.7748/0.7748 =
1.00$, el cociente de la \omterm{def:b1:weathering-clays:clay-mineral}{caolinita}, \omterm{def:b1:weathering-clays:hydrolysis-degree}{monosialitización}. B: $0.9987/0.4904 = 2.04$, una
arcilla 2:1, \omterm{def:b1:weathering-clays:hydrolysis-degree}{bisialitización}. C: $0.0333/1.2358 = 0.027$, casi no queda silicio: gibbsita, \omterm{def:b1:weathering-clays:hydrolysis-degree}{alitización}. De la lixiviación menos intensa a la más
intensa, es decir, hacia el clima más cálido y húmedo: B, A, C.
\end{solution}

\begin{solution}{pb:b1:weathering-clays:1}
\textbf{1.} \ce{2 NaAlSi3O8 + 2 H+ + 9 H2O -> Al2Si2O5(OH)4 + 2 Na+ + 4 H4SiO4}.

\textbf{2.} \ce{CaAl2Si2O8 + 2 H+ + H2O -> Al2Si2O5(OH)4 + Ca^2+}. No sale silicio:
la anortita ya tiene el silicio por aluminio de la \omterm{def:b1:weathering-clays:clay-mineral}{caolinita}.

\textbf{3.} Una gibbsita toma un aluminio, media anortita: duplicando,
\ce{CaAl2Si2O8 + 2 H+ + 6 H2O -> 2 Al(OH)3 + Ca^2+ + 2 H4SiO4}. Oxígeno:
$8 + 6 = 6 + 8$; hidrógeno: $2 + 12 = 6 + 8$.

\textbf{4.} \ce{5 NaAlSi3O8 + 4 H+ + 16 H2O -> 2 Na_{0.5}Al2(Si_{3.5}Al_{0.5})O10(OH)2 + 4 Na+ + 8 H4SiO4}; % ce-unbalanced-ok: decimal subscripts the checker cannot read; balance tested in tests/bachelor-1/test_weathering_clays.py
hidrógeno, $4 + 32 = 36$ a la izquierda y $2\times2 + 8\times4 = 36$ a la derecha, equilibrado.

\textbf{5.} Cada \ce{H+} viene con un \ce{HCO3-} de un \ce{CO2} y un \ce{H2O}:
\ce{2 NaAlSi3O8 + 2 CO2 + 11 H2O -> Al2Si2O5(OH)4 + 2 Na+ + 2 HCO3- + 4 H4SiO4}
y \ce{CaAl2Si2O8 + 2 CO2 + 3 H2O -> Al2Si2O5(OH)4 + Ca^2+ + 2 HCO3-}.

\textbf{6.} La albita a beidellita es \omterm{def:b1:weathering-clays:hydrolysis-degree}{bisialitización}; la albita y la anortita a \omterm{def:b1:weathering-clays:clay-mineral}{caolinita} son \omterm{def:b1:weathering-clays:hydrolysis-degree}{monosialitización}; la anortita a gibbsita es \omterm{def:b1:weathering-clays:hydrolysis-degree}{alitización}.

\textbf{7.} $\num{e6}/262.22 = \qty{3813.6}{mol}$ de albita;
$\num{e6}/278.20 = \qty{3594.5}{mol}$ de anortita.

\textbf{8.} Albita: $3813.6/2 = \qty{1906.8}{mol}$ de \omterm{def:b1:weathering-clays:clay-mineral}{caolinita},
\qty{492.3}{kg}. Anortita: \qty{3594.5}{mol}, \qty{928.0}{kg}.

\textbf{9.} Albita: $2\times3813.6 = \qty{7627.2}{mol}$ de \ce{SiO2},
\qty{458.2}{kg}, y \qty{3813.6}{mol} de sodio, \qty{87.7}{kg}. Anortita:
ninguna sílice, y \qty{3594.5}{mol} de calcio, \qty{144.1}{kg}.

\textbf{10.} Albita: $1/2.62 = \qty{0.382}{m^3}$; su \omterm{def:b1:weathering-clays:clay-mineral}{caolinita},
$0.4923/2.6 = \qty{0.189}{m^3}$: el sólido pierde el \qty{50.4}{\percent} de
su volumen. Anortita: $1/2.75 = \qty{0.364}{m^3}$; su \omterm{def:b1:weathering-clays:clay-mineral}{caolinita},
$0.9280/2.6 = \qty{0.357}{m^3}$: una pérdida de solo el \qty{1.8}{\percent}.

\textbf{11.} El espacio que queda se convierte en poros, llenos de agua y aire: una
saprolita sobre un granito rico en albita es ligera, porosa y desmenuzable, con sus granos de cuarzo sueltos en un poco de arcilla. La anortita se convierte en casi
su propio volumen de \omterm{def:b1:weathering-clays:clay-mineral}{caolinita}, así que el gabro meteorizado conserva la mayor parte de su materia sólida.

\textbf{12.} El sodio ($z/r = 0.98$) y el calcio (2.0) son cationes solubles y salen; el aluminio (5.6) es un hidrolizado y queda por completo en la
\omterm{def:b1:weathering-clays:clay-mineral}{caolinita}; el silicio (10.0), en el límite, sale en parte de la albita y no sale de la anortita, cuyo silicio puede retener la \omterm{def:b1:weathering-clays:clay-mineral}{caolinita}.

\textbf{13.} Un \ce{CO2} por albita: \qty{3813.6}{mol}, $3813.6 \times 44.01 =
\qty{167.8}{kg}$.

\textbf{14.} Dos por anortita: \qty{7189.0}{mol}, \qty{316.4}{kg}.

\textbf{15.} $7189.0 \times 61.02 = \qty{438.7}{kg}$ de bicarbonato.

\textbf{16.} Los organismos lo precipitan como carbonato de calcio,
\ce{Ca^2+ + 2 HCO3- -> CaCO3 + CO2 + H2O}: por cada calcio, uno de los dos
\ce{CO2} vuelve al aire, \qty{3594.5}{mol} o \qty{158.2}{kg}, y el otro queda enterrado en el carbonato.

\textbf{17.} El sodio no forma carbonato en el fondo marino: su bicarbonato no puede almacenarse como roca carbonatada, y el almacenamiento a largo plazo se cuenta solo a partir del calcio y del magnesio.

\textbf{18.} Una tonelada de calcita consume \qty{9991}{mol} de \ce{CO2} al disolverse y los devuelve todos al precipitar: no se almacena nada, y solo la
meteorización de los silicatos de calcio y de magnesio retira dióxido de carbono para siempre.

\textbf{19.} $\qty{200}{m^3.s^{-1}} \times \qty{0.3}{mol.m^{-3}} =
\qty{60}{mol.s^{-1}}$; en un año de \qty{3.156e7}{s}, \qty{1.893e9}{mol}, o
$\num{1.893e9}/\num{50000} = \qty{37870}{mol}$ por kilómetro cuadrado.

\textbf{20.} Dos bicarbonatos por anortita: \qty{9.467e8}{mol} al año,
$\num{9.467e8}\times\qty{278.20}{g} = \qty{263400}{t}$ al año, o
\qty{5.27}{t} por kilómetro cuadrado.

\textbf{21.} La roca fresca consumida es $5.27/0.20 = \qty{26.3}{t}$, o
$\num{26340}/2700 = \qty{9.75}{m^3}$, por kilómetro cuadrado y año: una capa de
\qty{9.75}{\micro m} de espesor al año, \qty{9.8}{m} por millón de años.

\textbf{22.} Un \ce{CO2} almacenado por calcio: $\num{9.467e8} \times 44.01 =
\qty{4.17e10}{g}$, unas \qty{41700}{t} al año.

\textbf{23.} Solo cuenta la materia arrastrada en disolución. La \omterm{def:b1:weathering-clays:clay-mineral}{caolinita} y el cuarzo que deja la meteorización se retiran también por la erosión,
como lodo y arena en el río, y no se cuentan; tampoco la meteorización de los demás \omterm{def:b1:rock-forming-minerals:mineral}{minerales} de la roca. Y si parte del bicarbonato
procede de los \omterm{def:b1:rock-forming-minerals:groups}{feldespatos} de sodio o de potasio, que dan un bicarbonato por unidad formular en lugar de dos, hizo falta meteorizar más roca para suministrarlo.

\textbf{24.} Según \cref{prop:b1:weathering-clays:co2}, se almacena un mol de \ce{CO2} por mol de anortita: $3594.5 \times \qty{44.01}{g} = \qty{158.19}{kg}$.
La meteorización de una tonelada de anortita retira para siempre \textbf{158}~kg de dióxido de carbono.
\end{solution}
